\begin{conclusion}[Lee--Yang 压缩有理函数的 Chebyshev 共轭与全显式反演]\label{con:xi-terminal-zm-leyang-compression-chebyshev-conjugacy-explicit}
沿用结论 \ref{con:xi-terminal-zm-leyang-discriminant-ridge-diff-square-norm} 的 \(A(y),Q(y)\)，并记
\[
a(y):=\frac{Q(y)}{8A(y)^3}\in\QQ(y).
\]
令 Chebyshev 三次多项式
\[
T_{3}(s):=s^{3}-3s.
\]
则存在 \(M\in\mathrm{PGL}_{2}(\QQ(\sqrt7))\) 与 \(N\in\mathrm{PGL}_{2}(\QQ(\sqrt7))\)（其中 \(N\) 可取为仿射变换）使得在 \(\QQ(\sqrt7)(y)\) 中成立严格恒等式
\[
N\bigl(a(y)\bigr)=T_{3}\bigl(M(y)\bigr).
\]
并且可以取为完全显式的
\[
M(y)=\frac{\sqrt7\,(8y-5)}{7(2y+1)},
\qquad
N(t)=\frac{\sqrt7}{49}\,(556-288t).
\]
等价地，
\[
\frac{\sqrt7}{49}\bigl(556-288a(y)\bigr)=\left(\frac{\sqrt7\,(8y-5)}{7(2y+1)}\right)^{3}-3\left(\frac{\sqrt7\,(8y-5)}{7(2y+1)}\right).
\]
从而 \(a:\PP^{1}_{y}\to\PP^{1}_{a}\) 在 \(\QQ(\sqrt7)\) 上与 \(T_{3}\) 线性分式共轭；其分歧型为 \((2,2,3)\)，与 \(T_{3}\) 完全一致。
\par
并且 \(M\) 可逆，其逆变换亦显式：
\[
M^{-1}(s)=\frac{-7s-5\sqrt7}{2(7s-4\sqrt7)}.
\]
因此对任意参数 \(t\)，若 \(s\) 为方程 \(s^{3}-3s=N(t)\) 的根，则 \(y=M^{-1}(s)\) 给出方程 \(a(y)=t\) 的一个根；这在 \(\QQ(\sqrt7)\) 上提供了 \(a\) 的全显式 Chebyshev 反演。
\end{conclusion}
\begin{proof}[证明要点]
由 \(a(y)=Q(y)/(8A(y)^3)\) 直接求导并化简可得
\[
a'(y)=-\frac{27}{8}\cdot\frac{2y^{2}-6y+1}{(2y+1)^{4}}.
\]
故 \(a\) 的有限临界点为
\(
y_{\pm}=\frac{3\pm\sqrt7}{2}
\)
（简单临界点，分歧指数 \(2\)），并且 \(y=-\tfrac12\) 为三阶极点，对应分歧指数 \(3\)。
将 \(y_{\pm}\) 代入 \(a(y)\) 得两条有限临界值
\[
a_{\pm}=\frac{139}{72}\pm\frac{7\sqrt7}{144},
\qquad
a\Bigl(-\frac12\Bigr)=\infty.
\]
取 \(M\) 为将三点 \((y_{-},y_{+},-\tfrac12)\) 送至 \((-1,1,\infty)\) 的唯一线性分式变换，即
\[
M(y)=\frac{\sqrt7\,(8y-5)}{7(2y+1)}.
\]
再取仿射 \(N\) 使得 \(N(a_{-})=2\)、\(N(a_{+})=-2\) 且 \(N(\infty)=\infty\)，解得
\[
N(t)=\frac{\sqrt7}{49}\,(556-288t).
\]
令
\(
R(y):=A(y)^{3}\bigl(N(a(y))-T_{3}(M(y))\bigr)
\)。
由于 \(N(a(y))\) 与 \(T_{3}(M(y))\) 的分母同为 \(A(y)^{3}\)，可知 \(R(y)\in\QQ(\sqrt7)[y]\) 且 \(\deg R\le 3\)。
由构造知 \(R(y_{\pm})=0\)。
又由结论 \ref{con:xi-terminal-zm-leyang-discriminant-ridge-diff-square-norm} 的恒等式
\(
a(y)-2=\frac{27y(y-1)}{8A(y)^{3}}
\)
得到 \(a(0)=a(1)=2\)，从而直接代入可核验 \(R(0)=R(1)=0\)。
于是 \(R\) 具有四个互异零点且 \(\deg R\le 3\)，故 \(R\equiv 0\)，即得所述共轭恒等式。
最后，\(M^{-1}\) 由线性分式求逆得到；Chebyshev 反演由 \(N(a(y))=T_{3}(M(y))\) 取根并代回 \(y=M^{-1}(s)\) 立即推出。
可复现核验脚本：\texttt{scripts/exp\_fold\_zm\_leyang\_compression\_chebyshev\_conjugacy\_audit.py}。证毕。
\end{proof}

\begin{conclusion}[\texorpdfstring{$a(y)$}{a(y)} 的三次逆像方程判别式约化与 \texorpdfstring{$S_{3}$}{S3} 伽罗瓦闭包；常数二次分辨支]\label{con:xi-terminal-zm-leyang-compression-cubic-preimage-discriminant-s3}
沿用结论 \ref{con:xi-terminal-zm-leyang-discriminant-ridge-diff-square-norm} 的 \(A(y),Q(y)\) 与结论 \ref{con:xi-terminal-zm-leyang-compression-chebyshev-conjugacy-explicit} 的 \(a(y)=Q(y)/(8A(y)^3)\)。
令 \(t\) 为 \(\PP^1_{a}\) 上坐标，并考虑关于 \(y\) 的三次多项式
\[
F_{t}(y):=Q(y)-8t\,A(y)^{3}\in\QQ(t)[y].
\]
则方程 \(a(y)=t\) 与 \(F_t(y)=0\) 等价（在 \(t\neq\infty\) 的仿射开集上）。
并且：
\begin{enumerate}
  \item \(F_t\) 在 \(\QQ(t)\) 上不可约。

  \item 其关于 \(y\) 的判别式满足严格因式分解
  \[
  \boxed{\;
  \mathrm{Disc}_{y}(F_t)=-3^{9}\,\bigl(2304t^{2}-8896t+8549\bigr)
  =-3^{9}\cdot 2304\,(t-t_{-})(t-t_{+})\; },
  \]
  其中
  \[
  t_{\pm}=\frac{139}{72}\pm\frac{7\sqrt7}{144}.
  \]

  \item 因为 \(\mathrm{Disc}_{y}(F_t)\notin(\QQ(t)^{\times})^{2}\)，故 \(F_t\) 的分裂域在 \(\QQ(t)\) 上的伽罗瓦群为 \(S_{3}\)。
  其唯一二次子扩张（分辨子域）为
  \[
  \QQ(t)\Bigl(\sqrt{2304t^{2}-8896t+8549}\Bigr).
  \]
  该二次分辨支的两条分歧点 \(t_{\pm}\) 生成常数二次域 \(\QQ(\sqrt7)\)，并且二次因子的判别式为 \(2^{10}\cdot 7^{3}\)。
\end{enumerate}
\end{conclusion}
\begin{proof}[证明要点]
不可约性：取特化 \(t=0\)，得到 \(F_{0}(y)=Q(y)\in\ZZ[y]\)。
由有理根定理可见 \(Q\) 无有理根，故在 \(\QQ\) 上不可约。
若 \(F_t\) 在 \(\QQ(t)\) 上可约，则除去有限多个极点后几乎所有有理特化皆可约，与 \(t=0\) 特化矛盾；故 \(F_t\) 在 \(\QQ(t)\) 上不可约。
\par
判别式分解：由结论 \ref{con:xi-terminal-zm-leyang-compression-chebyshev-conjugacy-explicit}，在 \(\QQ(\sqrt7)\) 上方程 \(a(y)=t\) 等价于
\[
T_{3}(s)=N(t),\qquad s=M(y),
\]
其中 \(T_{3}(s)-u=s^{3}-3s-u\) 的判别式为 \(\mathrm{Disc}_{s}(T_{3}(s)-u)=-27(u^{2}-4)\)。
线性分式代换 \(s=M(y)\) 与清分母仅改变判别式一个平方因子，故 \(\mathrm{Disc}_{y}(F_t)\) 与 \(-27\bigl(N(t)^{2}-4\bigr)\) 在 \(\QQ(\sqrt7)(t)\) 中相差一个平方因子。
由 \(N(t)=\frac{\sqrt7}{49}(556-288t)\) 直接计算得到严格恒等式
\[
N(t)^{2}-4=\frac{36}{343}\,\bigl(2304t^{2}-8896t+8549\bigr).
\]
另一方面，对 \(F_t\) 直接计算 \(\mathrm{Disc}_{y}(F_t)\) 并因式分解，即得到所示严格分解式与 \(t_{\pm}\)。
其中二次多项式的判别式为
\(
8896^{2}-4\cdot 2304\cdot 8549=2^{10}\cdot 7^{3}
\)，从而 \(t_{\pm}\in\QQ(\sqrt7)\setminus\QQ\)。
\par
伽罗瓦群：对不可约三次，多项式判别式（或其首项系数约化形式）为平方当且仅当伽罗瓦群为 \(A_{3}\)。
此处 \(\mathrm{Disc}_{y}(F_t)\) 在 \(\QQ(t)\) 中非平方，故伽罗瓦群为 \(S_{3}\)，其唯一二次子扩张由 \(\sqrt{\mathrm{Disc}_{y}(F_t)}\) 给出，即为所述分辨子域。
可复现核验脚本：\texttt{scripts/exp\_fold\_zm\_leyang\_compression\_chebyshev\_conjugacy\_audit.py}。证毕。
\end{proof}

\begin{conclusion}[Chebyshev 扭曲的 \texorpdfstring{$\QQ$}{Q}-不可消去性与最小共轭域]\label{con:xi-terminal-zm-leyang-compression-minimal-conjugacy-field}
设 \(a(y)\in\QQ(y)\) 如结论 \ref{con:xi-terminal-zm-leyang-compression-chebyshev-conjugacy-explicit}。
则 \(a\) 在 \(\QQ\) 上不可能与 \(T_{3}(s)=s^{3}-3s\) 线性分式共轭；
然而在 \(\QQ(\sqrt7)\) 上存在显式共轭，故 \(a\) 与 \(T_3\) 的最小共轭域恰为 \(\QQ(\sqrt7)\)。
\end{conclusion}
\begin{proof}[证明要点]
由结论 \ref{con:xi-terminal-zm-leyang-compression-chebyshev-conjugacy-explicit} 的导数计算，\(a\) 的两个有限临界点为
\(
y_{\pm}=\frac{3\pm\sqrt7}{2}
\notin\QQ
\)；
且其为共轭二次点。
设反证存在 \(\alpha,\beta\in\mathrm{PGL}_{2}(\QQ)\) 使 \(\beta\circ a\circ\alpha^{-1}=T_{3}\)。
后合 \(\beta\) 不改变临界点，故 \(T_{3}\) 的有限临界点集合 \(\{\pm 1\}\subset\QQ\) 必等于 \(\alpha(\{y_{+},y_{-}\})\)。
但由于 \(\alpha\) 的系数在 \(\QQ\) 中，\(\alpha(y_{+})\) 与 \(\alpha(y_{-})\) 仍为共轭二次数；
若其中之一为有理数，则另一亦为同一有理数，矛盾于 \(\alpha\) 的单射性。
故 \(\QQ\)-共轭不可能成立。
\par
另一方面，结论 \ref{con:xi-terminal-zm-leyang-compression-chebyshev-conjugacy-explicit} 已给出 \(\QQ(\sqrt7)\)-共轭的显式构造。
因此最小共轭域即为 \(\QQ(\sqrt7)\)。证毕。
\end{proof}

\begin{corollary}[最小共轭域的分歧素数支撑：\texorpdfstring{$\{2,7\}$}{\{2,7\}}]\label{cor:xi-terminal-zm-leyang-chebyshev-min-conjugacy-ramification-2-7}
在结论 \ref{con:xi-terminal-zm-leyang-compression-minimal-conjugacy-field} 的记号下，最小共轭域 \(K_7=\QQ(\sqrt7)\) 的基本判别式为 \(\Disc(K_7)=28\)，从而其分歧素数集合为
\[
\mathrm{Ram}(K_7)=\{2,7\}.
\]
因此在“局部提升一致 + 可逆实现”之记录轴口径下，任何实现该 Chebyshev 去扭共轭的素数账本支撑都必须包含 \(\{2,7\}\)，并且在二次数域情形下该约束为最小约束。
\end{corollary}
\begin{proof}
\(\Disc(\QQ(\sqrt7))=28\) 为标准计算，故分歧素数为其素因子集合 \(\{2,7\}\)。
最后的“最小性”由定理 \ref{thm:conclusion-quadratic-field-min-prime-ledger-discriminant}（二次数域情形）推出。证毕。
\end{proof}
\leanverified{paper\_xi\_terminal\_zm\_leyang\_chebyshev\_min\_conjugacy\_ramification\_2\_7}

\paragraph{压缩函数的 Chebyshev 线性化与平方下降生成元的判别式结构}\label{par:xi-terminal-zm-leyang-chebyshev-square-descent-index283}
固定记号
\[
Q(y):=16+69y+219y^2+128y^3,\qquad
A(y):=1+2y,\qquad
a(y):=\frac{Q(y)}{8A(y)^3},
\]
\[
P_{\mathrm{LY}}(y):=256y^3+411y^2+165y+32.
\]
并有恒等式
\[
a(y)-2=\frac{27y(y-1)}{8A(y)^3},\qquad
a(y)+2=\frac{P_{\mathrm{LY}}(y)}{8A(y)^3}.
\]
因此
\(
a^{-1}(2)=\{0,1,\infty\}
\)，且
\(
a^{-1}(-2)=\{P_{\mathrm{LY}}=0\}
\)（按重数计）。

\begin{proposition}[压缩函数 \texorpdfstring{$a(y)$}{a(y)} 的 Chebyshev 线性化（奇归一化）]\label{prop:xi-terminal-zm-leyang-chebyshev-linearization-odd-normalization}
\leanverified{paper\_xi\_terminal\_zm\_leyang\_chebyshev\_linearization\_odd\_normalization}
令
\[
\psi(s):=s^3-3s,
\qquad
\mu(y):=\frac{\sqrt7\,(5-8y)}{7(2y+1)},
\qquad
\nu(u):=\frac{288\sqrt7}{49}u-\frac{556\sqrt7}{49}.
\]
则在 \(\QQ(\sqrt7)(y)\) 中成立严格恒等式
\[
\nu\!\bigl(a(y)\bigr)=\psi\!\bigl(\mu(y)\bigr).
\]
\end{proposition}

\begin{proof}
由结论 \ref{con:xi-terminal-zm-leyang-compression-chebyshev-conjugacy-explicit}，
已知
\[
N\!\bigl(a(y)\bigr)=\psi\!\bigl(M(y)\bigr),\qquad
M(y)=\frac{\sqrt7\,(8y-5)}{7(2y+1)},\qquad
N(u)=\frac{\sqrt7}{49}(556-288u).
\]
注意 \(\mu=-M\)、\(\nu=-N\)，且 \(\psi\) 为奇函数，故
\[
\nu(a)= -N(a)= -\psi(M)=\psi(-M)=\psi(\mu).
\]
证毕。
\end{proof}

\begin{theorem}[分歧值域的二次指纹与最小共轭域]\label{thm:xi-terminal-zm-leyang-chebyshev-branch-quadratic-fingerprint}
映射 \(a:\PP^1_y\to\PP^1\) 的分歧点与分歧值满足
\[
\mathrm{Ram}(a)=\left\{-\frac12,\ \frac{3+\sqrt7}{2},\ \frac{3-\sqrt7}{2}\right\},
\]
其中分歧指数分别为 \(3,2,2\)，并且
\[
\mathrm{Br}(a)=\left\{\infty,\ \frac{139}{72}\pm\frac{7\sqrt7}{144}\right\}.
\]
此外，若存在域 \(L\) 及 \(\mu_0,\nu_0\in\mathrm{PGL}_2(L)\) 使
\(
\nu_0\circ a=\psi\circ\mu_0
\)，则必有 \(L\supseteq \QQ(\sqrt7)\)；
反之在 \(\QQ(\sqrt7)\) 上存在此类共轭。
故实现该 Chebyshev 线性化的最小共轭域为
\[
K_7=\QQ(\sqrt7).
\]
\end{theorem}

\begin{proof}
由
\[
a'(y)=-\frac{27\,(2y^2-6y+1)}{8(2y+1)^4}
\]
可得有限临界点 \(y_\pm=(3\pm\sqrt7)/2\)，且 \(y=-1/2\) 为三重极点，对应全分歧点。
代入 \(a(y)\) 得
\[
a(y_\pm)=\frac{139}{72}\pm\frac{7\sqrt7}{144}.
\]
从而得到 \(\mathrm{Ram}(a)\) 与 \(\mathrm{Br}(a)\)。

若 \(\nu_0\circ a=\psi\circ\mu_0\)，则
\(
\mathrm{Br}(a)=\nu_0^{-1}(\{\pm2,\infty\})\subset\PP^1(L)
\)。
而 \(\mathrm{Br}(a)\) 含两条互为共轭且非有理元素
\(
\frac{139}{72}\pm\frac{7\sqrt7}{144}
\)，故 \(L\supseteq\QQ(\sqrt7)\)。
反向存在性由命题 \ref{prop:xi-terminal-zm-leyang-chebyshev-linearization-odd-normalization} 给出。证毕。
\end{proof}
\leanverified{paper\_xi\_terminal\_zm\_leyang\_chebyshev\_branch\_quadratic\_fingerprint}

\begin{theorem}[Lee--Yang 三次根在 Chebyshev 坐标中的三次方程]\label{thm:xi-terminal-zm-leyang-ply-root-chebyshev-cubic}
令 \(\beta\) 为 \(P_{\mathrm{LY}}(y)=0\) 的任一根，并置
\[
s:=\mu(\beta)=\frac{\sqrt7\,(5-8\beta)}{7(2\beta+1)}.
\]
则 \(s\) 满足
\[
49s^3-147s+1132\sqrt7=0,
\]
等价地
\[
s^3-3s=-\frac{1132\sqrt7}{49}=\nu(-2).
\]
\end{theorem}

\begin{proof}
由 \(a(y)+2=P_{\mathrm{LY}}(y)/(8A(y)^3)\) 可知 \(a(\beta)=-2\)。
再由命题 \ref{prop:xi-terminal-zm-leyang-chebyshev-linearization-odd-normalization},
\[
\psi(s)=\psi(\mu(\beta))=\nu(a(\beta))=\nu(-2).
\]
直接计算
\[
\nu(-2)=\frac{288\sqrt7}{49}(-2)-\frac{556\sqrt7}{49}
=-\frac{1132\sqrt7}{49}.
\]
代入 \(\psi(s)=s^3-3s\) 即得结论。证毕。
\end{proof}
\leanverified{paper\_xi\_terminal\_zm\_leyang\_ply\_root\_chebyshev\_cubic}

\begin{proposition}[平方下降生成元与首一整式模型]\label{prop:xi-terminal-zm-leyang-square-descent-monic-generator}
令 \(\beta\) 为 \(P_{\mathrm{LY}}(y)=0\) 的任一根，定义
\[
T:=\left(\frac{5-8\beta}{2\beta+1}\right)^2\in\QQ(\beta).
\]
则 \(T\) 满足首一三次方程
\[
F(T):=T^3-42T^2+441T-1281424=0,
\]
并且
\[
\QQ(T)=\QQ(\beta).
\]
\end{proposition}

\begin{proof}
对不定元 \(y\) 记
\[
T(y):=\left(\frac{5-8y}{2y+1}\right)^2.
\]
计算结果式
\[
\mathrm{Res}_y\!\left(P_{\mathrm{LY}}(y),\ (5-8y)^2-T(2y+1)^2\right)
=26244\,F(T).
\]
故代入 \(y=\beta\) 得 \(F(T)=0\)。

由有理根判别可知 \(F\) 在 \(\QQ\) 上无有理根，故三次 \(F\) 不可约，从而
\(
[\QQ(T):\QQ]=3
\)。
又 \(T\in\QQ(\beta)\) 且 \([\QQ(\beta):\QQ]=3\)（见结论 \ref{con:xi-terminal-zm-leyang-cubic-field-isomorphism}），于是
\[
\QQ(T)=\QQ(\beta).
\]
证毕。
\end{proof}
\leanverified{paper\_xi\_terminal\_zm\_leyang\_square\_descent\_monic\_generator}

\begin{theorem}[新生成元的判别式与整指数]\label{thm:xi-terminal-zm-leyang-square-descent-discriminant-index}
沿用命题 \ref{prop:xi-terminal-zm-leyang-square-descent-monic-generator} 的 \(F(T)\)。
设 \(K=\QQ(\beta)\)，并记 \(\mathcal O_K\) 为其整数环。
则
\[
\mathrm{Disc}(F)= -2^6\cdot 3^5\cdot 31^2\cdot 37\cdot 283^2.
\]
若再使用结论 \ref{con:xi-terminal-zm-leyang-cubic-field-minimal-discriminant}
中的
\(
\mathrm{Disc}(K)=-999=-3^3\cdot 37
\)，则
\[
[\mathcal O_K:\ZZ[T]]
=2^3\cdot 3\cdot 31\cdot 283.
\]
\end{theorem}

\begin{proof}
对首一三次 \(F(T)\) 直接计算判别式并在 \(\ZZ\) 中分解，得到上式。
另一方面，
\[
\mathrm{Disc}(\ZZ[T])=\mathrm{Disc}(K)\,[\mathcal O_K:\ZZ[T]]^2.
\]
故
\[
[\mathcal O_K:\ZZ[T]]^2
=\frac{\mathrm{Disc}(F)}{\mathrm{Disc}(K)}
=2^6\cdot 3^2\cdot 31^2\cdot 283^2,
\]
从而指数为其正平方根
\(
2^3\cdot 3\cdot 31\cdot 283
\)。
证毕。
\end{proof}
\leanverified{paper\_xi\_terminal\_zm\_leyang\_square\_descent\_discriminant\_index}

\begin{remark}[可复现核验]\label{rem:xi-terminal-zm-leyang-chebyshev-square-descent-audit}
上述恒等式、结果式、判别式与指数关系可由脚本
\texttt{scripts/exp\_fold\_zm\_leyang\_compression\_chebyshev\_conjugacy\_audit.py}
直接复核。
\end{remark}

\endinput


\endinput
